\renewcommand{\thesection}{\arabic{section}}\renewcommand{\theHsection}{chapter.\arabic{section}}\setcounter{section}{13}\section{数项级数}
\noindent\begin{minipage}{\textwidth}\ProblemHeading{例题}{14.1（改写）}
设 $q\in\mathbb R$。讨论等比级数 $\sum_{n=0}^{\infty}q^n$ 的敛散性，并在收敛时求其和。
\end{minipage}\par\vfill
\noindent\begin{minipage}{\textwidth}\ProblemHeading{例题}{14.2}
设 $p\ge0,q>0,s=p+q$。计算级数 $\displaystyle\sum_{n=1}^{\infty}\frac1{(sn-p)(sn+q)}$ 的和。
\end{minipage}\par\vfill
\noindent\begin{minipage}{\textwidth}\ProblemHeading{例题}{14.3（改写）}
证明级数 $\sum_{n=0}^{\infty}(-1)^n$ 发散。
\end{minipage}\par\vfill
\newpage
\noindent\begin{minipage}{\textwidth}\ProblemHeading{例题}{14.4（改写）}
证明级数 $\sum_{n=1}^{\infty}n\sin(1/n)$ 发散。
\end{minipage}\par\vfill
\noindent\begin{minipage}{\textwidth}\ProblemHeading{例题}{14.5}
求级数 $\displaystyle\sum_{n=0}^{\infty}\frac{(-1)^n+2^n}{3^{n+2}}$ 的和。
\end{minipage}\par\vfill
\noindent\begin{minipage}{\textwidth}\ProblemHeading{例题}{14.6（改写）}
将级数 $1-1+1-1+\cdots$ 的相邻两项依次配成一组，比较配组前后两个级数的敛散性。据此判断：在不改变项的次序时，配组所得级数收敛，能否推出原级数收敛？
\end{minipage}\par\vfill
\newpage
\noindent\begin{minipage}{\textwidth}\ProblemHeading{习题}{14.7}
判断下列级数的敛散性。
\SubQ{1}{$\displaystyle\sum_{n=1}^{\infty}n$。}
\SubQ{2}{$\displaystyle\sum_{n=1}^{\infty}\frac1{\sqrt n}$。}
\SubQ{3}{$\displaystyle\sum_{n=0}^{\infty}\frac1{5^n}$。}
\end{minipage}\par\vfill
\noindent\begin{minipage}{\textwidth}\ProblemHeading{习题}{14.8}
判断下列级数的敛散性。
\SubQ{1}{$\displaystyle\sum_{n=0}^{\infty}\frac1{2n+3}$。}
\SubQ{2}{$\displaystyle\sum_{n=2}^{\infty}\frac1{\ln n}$。}
\end{minipage}\par\vfill
\noindent\begin{minipage}{\textwidth}\ProblemHeading{习题}{14.9}
求下列级数的和。
\SubQ{1}{$\displaystyle\sum_{n=1}^{\infty}\frac{(-1)^{n-1}}{2^{n-1}}$。}
\SubQ{2}{$\displaystyle\sum_{n=1}^{\infty}\frac1{(2n-1)(2n+1)}$。}
\SubQ{3}{$\displaystyle\sum_{n=1}^{\infty}\frac1{(3n-2)(3n+1)}$。}
\end{minipage}\par\vfill
\newpage
\noindent\begin{minipage}{\textwidth}\ProblemHeading{习题}{14.10}
求下列级数的和。
\SubQ{1}{$\displaystyle\sum_{n=1}^{\infty}\frac1{n(n+3)}$。}
\SubQ{2}{$\displaystyle\sum_{n=1}^{\infty}\frac{2n-1}{2^n}$。}
\end{minipage}\par\vfill
\noindent\begin{minipage}{\textwidth}\ProblemHeading{习题}{14.11}
证明下列等式。
\SubQ{1}{$\displaystyle\sum_{n=1}^{\infty}\frac{2n+1}{n^2(n+1)^2}=1$。}
\SubQ{2}{$\displaystyle\sum_{n=1}^{\infty}(\sqrt{n+2}-2\sqrt{n+1}+\sqrt n)=1-\sqrt2$。}
\SubQ{3}{$\displaystyle\sum_{n=1}^{\infty}\ln\frac{n(2n+1)}{(n+1)(2n-1)}=\ln2$。}
\end{minipage}\par\vfill
\noindent\begin{minipage}{\textwidth}\ProblemHeading{习题}{14.12}
证明下列等式。
\SubQ{1}{$\displaystyle\sum_{n=1}^{\infty}\frac1{n(n+m)}=\frac1m\left(1+\frac12+\cdots+\frac1m\right)$，$m$ 为正整数。}
\end{minipage}\par\vfill
\newpage
\noindent\begin{minipage}{\textwidth}\ProblemHeading{习题}{14.13}
构造一个无穷级数，使其第 $n$ 个部分和为 $S_n=1/n$（$n\ge1$）。
\end{minipage}\par\vfill
\noindent\begin{minipage}{\textwidth}\ProblemHeading{习题}{14.14}
若级数 $\sum_{n=1}^{\infty}a_n$ 与 $\sum_{n=1}^{\infty}b_n$ 都发散，对下列级数的敛散性能得出什么结论？
\SubQ{1}{$\displaystyle\sum_{n=1}^{\infty}(a_n+b_n)$。}
\SubQ{2}{$\displaystyle\sum_{n=1}^{\infty}(a_n-b_n)$。}
\SubQ{3}{$\displaystyle\sum_{n=1}^{\infty}a_nb_n$。}
\end{minipage}\par\vfill
\noindent\begin{minipage}{\textwidth}\ProblemHeading{习题}{14.15}
若级数 $\sum_{n=1}^{\infty}a_n$ 与 $\sum_{n=1}^{\infty}b_n$ 都发散，对下列级数的敛散性能得出什么结论？
\SubQ{1}{$\displaystyle\sum_{n=1}^{\infty}\frac{a_n}{b_n}$，假设每个 $b_n\ne0$。}
\end{minipage}\par\vfill
\newpage
\noindent\begin{minipage}{\textwidth}\ProblemHeading{习题}{14.16}
证明下列级数发散。
\SubQ{1}{$\displaystyle\sum_{n=1}^{\infty}\frac n{n+1}$。}
\SubQ{2}{$\displaystyle\sum_{n=1}^{\infty}(-1)^n\frac{n^2+1}{3n^2-2}$。}
\SubQ{3}{$\displaystyle\sum_{n=1}^{\infty}\frac1{\sqrt[n]n}$。}
\end{minipage}\par\vfill
\noindent\begin{minipage}{\textwidth}\ProblemHeading{习题}{14.17}
证明下列级数发散。
\SubQ{1}{$\displaystyle\sum_{n=1}^{\infty}\left(1-\frac1n\right)^n$。}
\end{minipage}\par\vfill
\noindent\begin{minipage}{\textwidth}\ProblemHeading{习题}{14.18}
设 $\sum_{n=1}^{\infty}a_n$ 收敛。证明 $\sum_{n=1}^{\infty}(a_n+a_{n+1})$ 也收敛。举例说明其逆命题不成立；再证明当所有 $a_n>0$ 时逆命题成立。
\end{minipage}\par\vfill
\newpage
\noindent\begin{minipage}{\textwidth}\ProblemHeading{习题}{14.19}
设数列 $\{na_n\}$ 与级数 $\sum_{n=1}^{\infty}n(a_n-a_{n+1})$ 都收敛。证明 $\sum_{n=1}^{\infty}a_n$ 收敛。
\end{minipage}\par\vfill
\noindent\begin{minipage}{\textwidth}\ProblemHeading{习题}{14.20（校订）}
设 $r$ 为正整数，$p\ne0$，且 $pn+q\ne0$ 对所有正整数 $n$ 成立。证明\[\sum_{n=1}^{\infty}\frac1{(pn+q)(pn+q+pr)}=\frac1{pr}\sum_{k=1}^r\frac1{q+kp}.\]
\end{minipage}\par\vfill
\noindent\begin{minipage}{\textwidth}\ProblemHeading{习题}{14.21}
设 $m$ 为给定正整数。证明\[\sum_{\substack{n=1\\n\ne m}}^{\infty}\frac1{m^2-n^2}=-\frac3{4m^2}.\]
\end{minipage}\par\vfill
\newpage
\noindent\begin{minipage}{\textwidth}\ProblemHeading{习题}{14.22}
证明 $\displaystyle12<1+\frac12+\frac13+\cdots+\frac1{10^6}<15$。
\end{minipage}\par\vfill
\noindent\begin{minipage}{\textwidth}\ProblemHeading{习题}{14.23}
设 $a_n>0$，数列 $\{a_n-a_{n+1}\}$ 严格递减，且 $\sum_{n=1}^{\infty}a_n$ 收敛。证明\[\lim_{n\to\infty}\left(\frac1{a_{n+1}}-\frac1{a_n}\right)=+\infty.\]
\end{minipage}\par\vfill
\noindent\begin{minipage}{\textwidth}\ProblemHeading{习题}{14.24}
证明存在与数列无关的正常数 $K$，使任意正数列 $\{a_n\}$ 满足\[\sum_{n=1}^{\infty}\frac n{a_1+\cdots+a_n}\le K\sum_{n=1}^{\infty}\frac1{a_n}.\]
\end{minipage}\par\vfill
\newpage
\noindent\begin{minipage}{\textwidth}\ProblemHeading{例题}{14.25}
设 $a_n>0$，$S_n=a_1+\cdots+a_n$。证明 $\sum_{n=1}^{\infty}a_n/S_n^2<+\infty$。
\end{minipage}\par\vfill
\noindent\begin{minipage}{\textwidth}\ProblemHeading{例题}{14.26（改写）}
证明级数 $\displaystyle\sum_{n=1}^{\infty}\frac1{\sqrt{n(n^2+1)}}$ 收敛。
\end{minipage}\par\vfill
\noindent\begin{minipage}{\textwidth}\ProblemHeading{例题}{14.27（改写）}
证明级数 $\sum_{n=2}^{\infty}(\ln n)^{-\ln n}$ 收敛。
\end{minipage}\par\vfill
\newpage
\noindent\begin{minipage}{\textwidth}\ProblemHeading{例题}{14.28（改写）}
证明级数 $\displaystyle\sum_{n=1}^{\infty}\frac{n+1}{\sqrt{n^3+2n-1}}$ 发散。
\end{minipage}\par\vfill
\noindent\begin{minipage}{\textwidth}\ProblemHeading{例题}{14.29}
对实数 $x$，研究级数 $\sum_{n=1}^{\infty}(1-\cos(x/n))$ 的敛散性。
\end{minipage}\par\vfill
\noindent\begin{minipage}{\textwidth}\ProblemHeading{例题}{14.30（改写）}
用积分判别法讨论 $p$ 级数 $\sum_{n=1}^{\infty}1/n^\alpha$ 的敛散性，其中 $\alpha\in\mathbb R$。
\end{minipage}\par\vfill
\newpage
\noindent\begin{minipage}{\textwidth}\ProblemHeading{例题}{14.31}
讨论级数 $\displaystyle\sum_{n=2}^{\infty}\frac1{n(\ln n)^\alpha}$ 的敛散性，其中 $\alpha\in\mathbb R$。
\end{minipage}\par\vfill
\noindent\begin{minipage}{\textwidth}\ProblemHeading{习题}{14.32}
设 $a_n\le b_n$（$n\ge1$）。若 $\sum_{n=1}^{\infty}b_n$ 收敛，能否断言 $\sum_{n=1}^{\infty}a_n$ 收敛？举例说明。
\end{minipage}\par\vfill
\noindent\begin{minipage}{\textwidth}\ProblemHeading{习题}{14.33}
用比较判别法讨论下列级数的敛散性。
\SubQ{1}{$\displaystyle\sum_{n=1}^{\infty}\frac1{3n^2+5}$。}
\SubQ{2}{$\displaystyle\sum_{n=1}^{\infty}\frac1{n2^n}$。}
\SubQ{3}{$\displaystyle\sum_{n=1}^{\infty}\left(\frac{n^2}{3n^2+1}\right)^n$。}
\end{minipage}\par\vfill
\newpage
\noindent\begin{minipage}{\textwidth}\ProblemHeading{习题}{14.34}
用比较判别法讨论下列级数的敛散性。
\SubQ{1}{$\displaystyle\sum_{n=1}^{\infty}\frac1n\sin\frac1n$。}
\SubQ{2}{$\displaystyle\sum_{n=1}^{\infty}\frac{n+1}{n(n+2)}$。}
\SubQ{3}{$\displaystyle\sum_{n=1}^{\infty}\frac1{n^{1+1/n}}$。}
\end{minipage}\par\vfill
\noindent\begin{minipage}{\textwidth}\ProblemHeading{习题}{14.35}
用比较判别法讨论下列级数的敛散性。
\SubQ{1}{$\displaystyle\sum_{n=3}^{\infty}\frac1{(\ln n)^{\ln\ln n}}$。}
\SubQ{2}{$\displaystyle\sum_{n=1}^{\infty}\left(n\ln\frac{2n+1}{2n-1}-1\right)$。}
\SubQ{3}{$\displaystyle\sum_{n=1}^{\infty}\left(n^{1/(n^2+1)}-1\right)$。}
\end{minipage}\par\vfill
\noindent\begin{minipage}{\textwidth}\ProblemHeading{习题}{14.36}
用比较判别法讨论下列级数的敛散性。
\SubQ{1}{$\displaystyle\sum_{n=1}^{\infty}\left(\frac1{\sqrt n}-\sqrt{\ln\frac{n+1}n}\right)$。}
\end{minipage}\par\vfill
\newpage
\noindent\begin{minipage}{\textwidth}\ProblemHeading{习题}{14.37}
若正项级数 $\sum_{n=1}^{\infty}a_n$ 收敛，证明 $\sum_{n=1}^{\infty}a_n^2$ 也收敛。举例说明逆命题不成立。
\end{minipage}\par\vfill
\noindent\begin{minipage}{\textwidth}\ProblemHeading{习题}{14.38}
若 $\sum_{n=1}^{\infty}a_n^2$ 与 $\sum_{n=1}^{\infty}b_n^2$ 收敛，证明 $\sum_{n=1}^{\infty}|a_nb_n|$ 与 $\sum_{n=1}^{\infty}(a_n+b_n)^2$ 也收敛。
\end{minipage}\par\vfill
\noindent\begin{minipage}{\textwidth}\ProblemHeading{习题}{14.39}
若正项级数 $\sum_{n=1}^{\infty}a_n$ 收敛，证明 $\sum_{n=1}^{\infty}\sqrt{a_na_{n+1}}$ 也收敛。举例说明逆命题不成立；再证明当 $\{a_n\}$ 递减时逆命题成立。
\end{minipage}\par\vfill
\newpage
\noindent\begin{minipage}{\textwidth}\ProblemHeading{习题}{14.40}
证明级数 $\sum_{n=2}^{\infty}1/\ln(n!)$ 发散。
\end{minipage}\par\vfill
\noindent\begin{minipage}{\textwidth}\ProblemHeading{习题}{14.41}
证明级数 $\displaystyle\sum_{n=3}^{\infty}\frac1{n\ln n\ln\ln n}$ 发散。
\end{minipage}\par\vfill
\noindent\begin{minipage}{\textwidth}\ProblemHeading{习题}{14.42}
当实数 $p,q$ 取何值时，级数 $\displaystyle\sum_{n=3}^{\infty}\frac1{n(\ln n)^p(\ln\ln n)^q}$ 收敛？
\end{minipage}\par\vfill
\newpage
\noindent\begin{minipage}{\textwidth}\ProblemHeading{习题}{14.43}
设 $\sum_{n=1}^{\infty}a_n$ 为收敛正项级数。证明对任意 $\delta>0$，有 $\sum_{n=1}^{\infty}n^{-(1+\delta)/2}\sqrt{a_n}<+\infty$。当 $\delta=0$ 时情况如何？
\end{minipage}\par\vfill
\noindent\begin{minipage}{\textwidth}\ProblemHeading{习题}{14.44}
设 $\sigma>0$，$a_n>0$。证明下列结论。
\SubQ{1}{若充分大的 $n$ 满足 $\ln(1/a_n)/\ln n\ge1+\sigma$，则 $\sum_{n=1}^{\infty}a_n$ 收敛。}
\SubQ{2}{若充分大的 $n$ 满足 $\ln(1/a_n)/\ln n\le1$，则 $\sum_{n=1}^{\infty}a_n$ 发散。}
\SubQ{3}{利用上述结论证明 $\sum_{n=2}^{\infty}1/(\ln n)^{\ln n}$ 与 $\sum_{n=1}^{\infty}1/3^{\ln n}$ 收敛。}
\end{minipage}\par\vfill
\noindent\begin{minipage}{\textwidth}\ProblemHeading{习题}{14.45}
设 $\sum_{n=1}^{\infty}a_n$ 为发散正项级数。证明 $\sum_{n=1}^{\infty}a_n/(1+a_n)$ 发散，而 $\sum_{n=1}^{\infty}a_n/(1+n^2a_n)$ 收敛。
\end{minipage}\par\vfill
\newpage
\noindent\begin{minipage}{\textwidth}\ProblemHeading{习题}{14.46}
设 $a_n>0$，$\sum_{n=1}^{\infty}a_n<+\infty$，$\alpha>0,\beta>0,\alpha+\beta>1$。证明 $\sum_{n=1}^{\infty}a_n^\alpha/n^\beta<+\infty$。
\end{minipage}\par\vfill
\noindent\begin{minipage}{\textwidth}\ProblemHeading{习题}{14.47}
证明级数 $\sum_{n=1}^{\infty}x^{1+1/2+\cdots+1/n}$ 在 $0<x<e^{-1}$ 时收敛，在 $x\ge e^{-1}$ 时发散。
\end{minipage}\par\vfill
\noindent\begin{minipage}{\textwidth}\ProblemHeading{习题}{14.48}
设 $\sum_{n=1}^{\infty}a_n$ 为发散正项级数。讨论 $\sum_{n=1}^{\infty}a_n/(1+a_n^2)$ 与 $\sum_{n=1}^{\infty}a_n/(1+na_n)$ 的敛散性。
\end{minipage}\par\vfill
\newpage
\noindent\begin{minipage}{\textwidth}\ProblemHeading{习题}{14.49}
设 $a_n>0$，$S_n=\sum_{k=1}^na_k$。证明：
\SubQ{1}{当 $\alpha>1$ 时，$\sum_{n=1}^{\infty}a_n/S_n^\alpha$ 收敛。}
\SubQ{2}{当 $\alpha\le1$ 且 $S_n\to+\infty$ 时，$\sum_{n=1}^{\infty}a_n/S_n^\alpha$ 发散。}
\end{minipage}\par\vfill
\noindent\begin{minipage}{\textwidth}\ProblemHeading{习题}{14.50}
设 $\{a_n\}$ 为递减正数列。证明下列结论。
\SubQ{1}{$\sum_{n=1}^{\infty}a_n$ 与 $\sum_{n=0}^{\infty}2^na_{2^n}$ 同敛散。}
\SubQ{2}{利用这一结论证明 $\sum_{n=1}^{\infty}1/n^{1+\alpha}$ 在 $\alpha>0$ 时收敛，在 $\alpha\le0$ 时发散。}
\SubQ{3}{利用这一结论证明 $\sum_{n=2}^{\infty}1/(n\ln n)$ 发散。}
\end{minipage}\par\vfill
\noindent\begin{minipage}{\textwidth}\ProblemHeading{习题}{14.51}
从调和级数 $\sum_{n=1}^{\infty}1/n$ 中去掉所有分母的十进制表示中含数字 $9$ 的项（如 $1/9,1/19,1/391$）。证明所得新级数收敛，且其和不超过 $80$。
\end{minipage}\par\vfill
\newpage
\noindent\begin{minipage}{\textwidth}\ProblemHeading{例题}{14.52（改写）}
分别计算级数 $\displaystyle\sum_{n=1}^{\infty}\frac1n$ 与 $\displaystyle\sum_{n=1}^{\infty}\frac1{n^2}$ 的通项的 $n$ 次方根极限，并比较两级数的敛散性。由此说明 Cauchy 根值判别法在极限等于 $1$ 时能否给出统一结论。
\end{minipage}\par\vfill
\noindent\begin{minipage}{\textwidth}\ProblemHeading{例题}{14.53}
讨论级数 $\displaystyle\sum_{n=1}^{\infty}\frac1{2^n}\left(1+\frac1n\right)^{n^2}$ 的敛散性。
\end{minipage}\par\vfill
\noindent\begin{minipage}{\textwidth}\ProblemHeading{例题}{14.54（改写）}
设正数列 $a_{2n-1}=2^{-n}$、$a_{2n}=3^{-n}$（$n\ge1$）。
\SubQ{1}{求 $\displaystyle\limsup_{n\to\infty}\sqrt[n]{a_n}$，并用 Cauchy 根值判别法讨论 $\displaystyle\sum_{n=1}^{\infty}a_n$ 的敛散性。}
\SubQ{2}{求 $\displaystyle\liminf_{n\to\infty}\frac{a_{n+1}}{a_n}$ 及 $\displaystyle\limsup_{n\to\infty}\frac{a_{n+1}}{a_n}$。说明 D’Alembert 比值判别法能否判断该级数的敛散性。}
\end{minipage}\par\vfill
\newpage
\noindent\begin{minipage}{\textwidth}\ProblemHeading{例题}{14.55}
设 $x\ge0$，讨论级数 $\sum_{n=1}^{\infty}n!x^n$ 的敛散性。
\end{minipage}\par\vfill
\noindent\begin{minipage}{\textwidth}\ProblemHeading{例题}{14.56}
设 $x\ge0$，讨论级数 $\sum_{n=1}^{\infty}(n!/n^n)x^n$ 的敛散性。
\end{minipage}\par\vfill
\noindent\begin{minipage}{\textwidth}\ProblemHeading{例题}{14.57（改写）}
设 $0<q<1$、$\alpha>1$、$\beta>1$。对下面各对级数，证明它们都收敛，且前一条级数的通项与后一条级数的通项之比趋于零。
\SubQ{1}{$\displaystyle\sum_{n=2}^{\infty}q^n$ 与 $\displaystyle\sum_{n=2}^{\infty}\frac1{n^\alpha}$。}
\SubQ{2}{$\displaystyle\sum_{n=2}^{\infty}\frac1{n^\alpha}$ 与 $\displaystyle\sum_{n=2}^{\infty}\frac1{n(\ln n)^\beta}$。}
\end{minipage}\par\vfill
\newpage
\noindent\begin{minipage}{\textwidth}\ProblemHeading{例题}{14.58}
设 $\alpha>0$，$\binom\alpha0=1$，$\binom\alpha n=\alpha(\alpha-1)\cdots(\alpha-n+1)/n!$。证明 $\sum_{n=0}^{\infty}|\binom\alpha n|$ 收敛。
\end{minipage}\par\vfill
\noindent\begin{minipage}{\textwidth}\ProblemHeading{例题}{14.59}
设 $p,q>0$。当 $p,q$ 取何值时，级数\[\sum_{n=1}^{\infty}\frac{p(p+1)\cdots(p+n-1)}{n!\,n^q}\]收敛？
\end{minipage}\par\vfill
\noindent\begin{minipage}{\textwidth}\ProblemHeading{例题}{14.60}
设 $\alpha,\beta,\gamma,x>0$。研究超几何级数\[F(\alpha,\beta,\gamma,x)=1+\sum_{n=1}^{\infty}\frac{\alpha(\alpha+1)\cdots(\alpha+n-1)\,\beta(\beta+1)\cdots(\beta+n-1)}{n!\,\gamma(\gamma+1)\cdots(\gamma+n-1)}x^n\]的敛散性。
\end{minipage}\par\vfill
\newpage
\noindent\begin{minipage}{\textwidth}\ProblemHeading{例题}{14.61（改写）}
设 $\alpha,\beta>0$。利用 Gauss 判别法研究级数\[1+\sum_{n=1}^{\infty}\frac{\alpha(\alpha+1)\cdots(\alpha+n-1)\,\beta(\beta+1)\cdots(\beta+n-1)}{n!\,(\alpha+\beta)(\alpha+\beta+1)\cdots(\alpha+\beta+n-1)}\]的敛散性。
\end{minipage}\par\vfill
\noindent\begin{minipage}{\textwidth}\ProblemHeading{习题}{14.62}
讨论下列级数的敛散性。
\SubQ{1}{$\displaystyle\sum_{n=1}^{\infty}n\tan\frac\pi{2^{n+1}}$。}
\SubQ{2}{$\displaystyle\sum_{n=1}^{\infty}\frac{n^2}{3^n}$。}
\SubQ{3}{$\displaystyle\sum_{n=1}^{\infty}\frac{n^5}{3^n}(\sqrt3+(-1)^n)^n$。}
\end{minipage}\par\vfill
\noindent\begin{minipage}{\textwidth}\ProblemHeading{习题}{14.63}
讨论下列级数的敛散性。
\SubQ{1}{$\displaystyle\sum_{n=1}^{\infty}\frac{n^2}{(1+1/n)^n}$。}
\SubQ{2}{$\displaystyle\sum_{n=1}^{\infty}\frac{n^{n+1/n}}{(n+1/n)^n}$。}
\SubQ{3}{$\displaystyle\sum_{n=1}^{\infty}\frac{x^n}{n!}$，$x\ge0$。}
\end{minipage}\par\vfill
\newpage
\noindent\begin{minipage}{\textwidth}\ProblemHeading{习题}{14.64}
讨论下列级数的敛散性。
\SubQ{1}{$\displaystyle\sum_{n=1}^{\infty}\left(\frac{n-4}{3n+1}\right)^n$。}
\SubQ{2}{$\displaystyle\sum_{n=2}^{\infty}\frac{n^{\ln n}}{(\ln n)^n}$。}
\end{minipage}\par\vfill
\noindent\begin{minipage}{\textwidth}\ProblemHeading{习题}{14.65}
利用 Raabe 判别法讨论下列级数的敛散性。
\SubQ{1}{$\displaystyle\sum_{n=1}^{\infty}\frac{\sqrt{n!}}{(a+\sqrt1)(a+\sqrt2)\cdots(a+\sqrt n)}$，$a>0$。}
\SubQ{2}{$\displaystyle\sum_{n=1}^{\infty}\frac{n!n^{-p}}{q(q+1)\cdots(q+n)}$，$p,q>0$。}
\end{minipage}\par\vfill
\noindent\begin{minipage}{\textwidth}\ProblemHeading{习题}{14.66}
设 $\{a_n\}$ 为正数列。证明\[\liminf_{n\to\infty}\frac{a_{n+1}}{a_n}\le\liminf_{n\to\infty}\sqrt[n]{a_n}.\]
\end{minipage}\par\vfill
\newpage
\noindent\begin{minipage}{\textwidth}\ProblemHeading{习题}{14.67}
设 $a_n>0$，且 $\lim_{n\to\infty}n(a_n/a_{n+1}-1)=l$。证明：$l>1$ 时 $\sum a_n$ 收敛，$l<1$ 时发散；并说明 $l=1$ 时不能据此判断敛散性。
\end{minipage}\par\vfill
\noindent\begin{minipage}{\textwidth}\ProblemHeading{习题}{14.68}
设 $p>0$。利用 Gauss 判别法判断级数\[\sum_{n=1}^{\infty}\frac{p(p+1)\cdots(p+n-1)}{n!\,n^p}\]的敛散性。
\end{minipage}\par\vfill
\noindent\begin{minipage}{\textwidth}\ProblemHeading{习题}{14.69}
设 $a_n>0$，存在常数 $\alpha,\beta>0$，使 $a_n-a_{n+1}\ge\beta a_n^{2-\alpha}$。证明 $\sum_{n=1}^{\infty}a_n<+\infty$，而且\[\sum_{n=N}^{\infty}a_n=O(a_N^\alpha)\quad(N\to\infty).\]
\end{minipage}\par\vfill
\newpage
\noindent\begin{minipage}{\textwidth}\ProblemHeading{习题}{14.70}
设 $a_n>0$，且 $\limsup_{n\to\infty}a_{n+1}/a_n=q<1$ 或 $\limsup_{n\to\infty}\sqrt[n]{a_n}=q<1$。证明对任意 $r\in(q,1)$，有 $\lim_{n\to\infty}a_n/r^n=0$。
\end{minipage}\par\vfill
\noindent\begin{minipage}{\textwidth}\ProblemHeading{习题}{14.71}
设 $\sum_{n=1}^{\infty}a_n$ 为收敛正项级数。构造一个收敛正项级数 $\sum_{n=1}^{\infty}b_n$，使 $\lim_{n\to\infty}a_n/b_n=0$。
\end{minipage}\par\vfill
\noindent\begin{minipage}{\textwidth}\ProblemHeading{例题}{14.72}
设 $\{a_n\}$ 为递减正数列。若 $\sum_{n=1}^{\infty}a_n$ 收敛，证明 $\lim_{n\to\infty}na_n=0$。
\end{minipage}\par\vfill
\newpage
\noindent\begin{minipage}{\textwidth}\ProblemHeading{例题}{14.73（改写）}
令 $a_n=1/(n\ln n)$（$n\ge2$）。验证 $\{a_n\}$ 递减且 $na_n\to0$，并讨论 $\displaystyle\sum_{n=2}^{\infty}a_n$ 的敛散性。由此判断递减正数列满足 $na_n\to0$ 时，其级数是否必收敛。
\end{minipage}\par\vfill
\noindent\begin{minipage}{\textwidth}\ProblemHeading{例题}{14.74（改写）}
分别用 Leibniz 判别法研究下列级数的敛散性，并判断收敛者属于绝对收敛还是条件收敛。
\SubQ{1}{$\displaystyle\sum_{n=1}^{\infty}\frac{(-1)^{n-1}}n$。}
\SubQ{2}{$\displaystyle\sum_{n=2}^{\infty}\frac{(-1)^n}{\ln n}$。}
\end{minipage}\par\vfill
\noindent\begin{minipage}{\textwidth}\ProblemHeading{例题}{14.75}
对实数 $x$，研究级数 $\sum_{n=1}^{\infty}\cos(nx)/n$ 的敛散性。
\end{minipage}\par\vfill
\newpage
\noindent\begin{minipage}{\textwidth}\ProblemHeading{例题}{14.76（改写）}
对 $x\in\mathbb R\setminus2\pi\mathbb Z$ 和正整数 $N$，求 $\displaystyle\sum_{n=1}^N\sin(nx)$ 的闭式，并证明其绝对值不超过 $1/|\sin(x/2)|$。
\end{minipage}\par\vfill
\noindent\begin{minipage}{\textwidth}\ProblemHeading{例题}{14.77}
研究级数 $\displaystyle\sum_{n=1}^{\infty}\frac{\cos3n}{n}\left(1+\frac1n\right)^n$ 的敛散性。
\end{minipage}\par\vfill
\noindent\begin{minipage}{\textwidth}\ProblemHeading{例题}{14.78}
研究级数 $\sum_{n=1}^{\infty}(-1)^{n-1}\sin^2n/n$ 的敛散性。
\end{minipage}\par\vfill
\newpage
\noindent\begin{minipage}{\textwidth}\ProblemHeading{习题}{14.79}
利用 Cauchy 收敛原理讨论下列级数的敛散性。
\SubQ{1}{$\displaystyle\sum_{n=1}^{\infty}\frac1{2n-1}$。}
\SubQ{2}{$\displaystyle\sum_{n=1}^{\infty}\frac{\sin n}{2^n}$。}
\SubQ{3}{$\displaystyle\sum_{n=1}^{\infty}\frac{\cos(n!)}{n(n+1)}$。}
\end{minipage}\par\vfill
\noindent\begin{minipage}{\textwidth}\ProblemHeading{习题}{14.80}
利用 Cauchy 收敛原理讨论下列级数的敛散性。
\SubQ{1}{$\displaystyle\sum_{n=1}^{\infty}\frac{a\cos n+b\sin n}{n\bigl(n+\sin(n!)\bigr)}$，$a,b\in\mathbb R$。}
\end{minipage}\par\vfill
\noindent\begin{minipage}{\textwidth}\ProblemHeading{习题}{14.81}
设 $a_n$ 递减趋于 $0$。利用 Cauchy 收敛原理证明 Leibniz 判别法：$\sum_{n=1}^{\infty}(-1)^{n-1}a_n$ 收敛。
\end{minipage}\par\vfill
\newpage
\noindent\begin{minipage}{\textwidth}\ProblemHeading{习题}{14.82}
设 $a_n\le c_n\le b_n$（$n\ge1$）。若 $\sum a_n$ 与 $\sum b_n$ 都收敛，证明 $\sum c_n$ 收敛。若前两个级数都发散，结论如何？
\end{minipage}\par\vfill
\noindent\begin{minipage}{\textwidth}\ProblemHeading{习题}{14.83}
证明级数\[\sum_{n=3}^{\infty}\frac{(-1)^{n-1}}{\sqrt{\lfloor(n+1)/2\rfloor}+(-1)^n}\]发散，并说明这表明 Leibniz 判别法中不能直接删去正数列递减的条件；其中 $\lfloor t\rfloor$ 表示整数部分。
\end{minipage}\par\vfill
\noindent\begin{minipage}{\textwidth}\ProblemHeading{习题}{14.84}
讨论下列级数的敛散性。
\SubQ{1}{$\displaystyle\sum_{n=1}^{\infty}(-1)^{n-1}\frac{\sqrt n}{n+1}$。}
\SubQ{2}{$\displaystyle\sum_{n=1}^{\infty}(-1)^{n-1}\frac1{\sqrt[n]n}$。}
\SubQ{3}{$\displaystyle\sum_{n=1}^{\infty}(-1)^{n-1}\sin\frac1n$。}
\end{minipage}\par\vfill
\newpage
\noindent\begin{minipage}{\textwidth}\ProblemHeading{习题}{14.85}
讨论下列级数的敛散性。
\SubQ{1}{$\displaystyle\sum_{n=1}^{\infty}\frac{\sin nx}n$，$x\in\mathbb R$。}
\end{minipage}\par\vfill
\noindent\begin{minipage}{\textwidth}\ProblemHeading{习题}{14.86}
设 $\{a_n\}$ 递减趋于 $0$。对实数 $x$，讨论 $\sum_{n=1}^{\infty}a_n\cos nx$ 与 $\sum_{n=1}^{\infty}a_n\sin nx$ 的敛散性。
\end{minipage}\par\vfill
\noindent\begin{minipage}{\textwidth}\ProblemHeading{习题}{14.87}
若 $\sum a_n$ 收敛且 $b_n/a_n\to1$，能否断言 $\sum b_n$ 收敛？研究级数 $\sum_{n=1}^{\infty}(-1)^{n-1}/\sqrt n$ 与 $\sum_{n=1}^{\infty}((-1)^{n-1}/\sqrt n+1/n)$ 以说明结论。
\end{minipage}\par\vfill
\newpage
\noindent\begin{minipage}{\textwidth}\ProblemHeading{习题}{14.88}
设 $\sum_{n=1}^{\infty}a_n$ 收敛。证明 $\lim_{n\to\infty}(a_1+2a_2+\cdots+na_n)/n=0$。
\end{minipage}\par\vfill
\noindent\begin{minipage}{\textwidth}\ProblemHeading{习题}{14.89}
证明：若 $\sum_{n=1}^{\infty}a_n/n^\alpha$ 收敛，则对任意 $\beta>\alpha$，级数 $\sum_{n=1}^{\infty}a_n/n^\beta$ 也收敛。
\end{minipage}\par\vfill
\noindent\begin{minipage}{\textwidth}\ProblemHeading{习题}{14.90}
设 $a_n$ 递减趋于 $0$，$p$ 为固定正整数。证明级数\[a_1+\cdots+a_p-a_{p+1}-\cdots-a_{2p}+a_{2p+1}+\cdots+a_{3p}-\cdots\]收敛。
\end{minipage}\par\vfill
\newpage
\noindent\begin{minipage}{\textwidth}\ProblemHeading{习题}{14.91}
讨论下列级数的敛散性。
\SubQ{1}{$\displaystyle\sum_{n=1}^{\infty}\sin(\pi\sqrt{n^2+1})$。}
\SubQ{2}{$\displaystyle\sum_{n=1}^{\infty}\left(1+\frac12+\cdots+\frac1n\right)\frac{\sin nx}n$，$x\in\mathbb R$。}
\end{minipage}\par\vfill
\noindent\begin{minipage}{\textwidth}\ProblemHeading{习题}{14.92}
设 $a_n>0$ 且 $\lim_{n\to\infty}n(a_n/a_{n+1}-1)=\lambda>0$。证明交错级数 $\sum_{n=1}^{\infty}(-1)^{n-1}a_n$ 收敛。
\end{minipage}\par\vfill
\noindent\begin{minipage}{\textwidth}\ProblemHeading{习题}{14.93}
设 $\{a_n\}$ 为正数列，$0<\alpha<1$，且 $\lim_{n\to\infty}n^\alpha(a_n/a_{n+1}-1)$ 为正数或 $+\infty$。证明对任意正整数 $k$，有 $\sum_{n=1}^{\infty}n^ka_n<+\infty$。
\end{minipage}\par\vfill
\newpage
\noindent\begin{minipage}{\textwidth}\ProblemHeading{习题}{14.94}
设 $a,b>0$。利用 Raabe 判别法证明，当 $b-1>a$ 时，\[1+\sum_{n=1}^{\infty}\frac{a(a+1)\cdots(a+n-1)}{b(b+1)\cdots(b+n-1)}\]收敛，且其和为 $(b-1)/(b-a-1)$。当 $b-1\le a$ 时敛散性如何？
\end{minipage}\par\vfill
\noindent\begin{minipage}{\textwidth}\ProblemHeading{习题}{14.95}
设 $\{a_n\}$ 为递增正数列。证明 $\sum_{n=1}^{\infty}(a_{n+1}/a_n-1)$ 收敛的充要条件是 $\{a_n\}$ 有界。
\end{minipage}\par\vfill
\noindent\begin{minipage}{\textwidth}\ProblemHeading{习题}{14.96}
证明级数 $\sum_{n=1}^{\infty}(-1)^{\lfloor\sqrt n\rfloor}/n$ 收敛，其中 $\lfloor t\rfloor$ 为整数部分。
\end{minipage}\par\vfill
\newpage
\noindent\begin{minipage}{\textwidth}\ProblemHeading{习题}{14.97}
构造一个收敛级数 $\sum_{n=1}^{\infty}a_n$，使 $\sum_{n=1}^{\infty}a_n^3$ 发散。
\end{minipage}\par\vfill
\noindent\begin{minipage}{\textwidth}\ProblemHeading{习题}{14.98}
设 $b_n$ 递减趋于 $0$，且 $\sum_{n=1}^{\infty}a_nb_n$ 收敛。证明 $\lim_{n\to\infty}(a_1+\cdots+a_n)b_n=0$。
\end{minipage}\par\vfill
\noindent\begin{minipage}{\textwidth}\ProblemHeading{习题}{14.99}
设 $\{b_n\}$ 为递增趋于 $+\infty$ 的正数列，且 $\sum_{n=1}^{\infty}a_n$ 收敛。证明\[\lim_{n\to\infty}\frac{a_1b_1+\cdots+a_nb_n}{b_n}=0.\]
\end{minipage}\par\vfill
\newpage
\noindent\begin{minipage}{\textwidth}\ProblemHeading{习题}{14.100}
设 $\sum_{n=1}^{\infty}a_n$ 收敛，$b_n=(a_1+2a_2+\cdots+na_n)/(n(n+1))$。证明 $\sum_{n=1}^{\infty}b_n$ 收敛，且其和等于 $\sum_{n=1}^{\infty}a_n$。
\end{minipage}\par\vfill
\noindent\begin{minipage}{\textwidth}\ProblemHeading{例题}{14.101（改写）}
判断下列各级数的绝对收敛性。记 $\lfloor t\rfloor$ 为不大于 $t$ 的最大整数，$x\in\mathbb R$。
\SubQ{1}{$\displaystyle\sum_{n=1}^{\infty}\frac{(-1)^{\lfloor\sqrt[3]{n}\rfloor}}{n^2}$。}
\SubQ{2}{$\displaystyle\sum_{n=1}^{\infty}\frac{\cos(nx)}{n^2}$。}
\SubQ{3}{$\displaystyle\sum_{n=1}^{\infty}\frac{\sin(nx)}{2^n}$。}
\end{minipage}\par\vfill
\noindent\begin{minipage}{\textwidth}\ProblemHeading{例题}{14.102}
证明级数 $\sum_{n=1}^{\infty}(-1)^n\sin^2n/n$ 条件收敛。
\end{minipage}\par\vfill
\newpage
\noindent\begin{minipage}{\textwidth}\ProblemHeading{例题}{14.103}
将交错调和级数 $1-\frac12+\frac13-\frac14+\cdots$ 重排为依次取 $p$ 个正项、$q$ 个负项，如此反复（各类项内部保持原次序），其中 $p,q$ 为正整数。判断所得级数是否收敛；若收敛，求其和。
\end{minipage}\par\vfill
\noindent\begin{minipage}{\textwidth}\ProblemHeading{习题}{14.104}
判断下列级数是绝对收敛还是条件收敛，其中 $x\in\mathbb R$。
\SubQ{1}{$\displaystyle\sum_{n=1}^{\infty}\frac{\sin nx}{n^\alpha}$，$\alpha>1$。}
\SubQ{2}{$\displaystyle\sum_{n=1}^{\infty}\frac{(-1)^{n-1}}{2^n}\cos nx$。}
\SubQ{3}{$\displaystyle\sum_{n=2}^{\infty}\frac{(-1)^n}{n\ln n}$。}
\end{minipage}\par\vfill
\noindent\begin{minipage}{\textwidth}\ProblemHeading{习题}{14.105}
判断下列级数是绝对收敛还是条件收敛，其中 $x\in\mathbb R$。
\SubQ{1}{$\displaystyle\sum_{n=1}^{\infty}(-1)^{n(n-1)/2}\frac{n^{10}}{a^n}$，$a>1$。}
\SubQ{2}{$\displaystyle\sum_{n=1}^{\infty}(-1)^n\frac{n-1}{n+1}\frac1{\sqrt[3]n}$。}
\SubQ{3}{$\displaystyle\sum_{n=2}^{\infty}\frac{\sin(n\pi/4)}{\ln n}$。}
\end{minipage}\par\vfill
\newpage
\noindent\begin{minipage}{\textwidth}\ProblemHeading{习题}{14.106}
讨论下列级数的绝对收敛性和条件收敛性，其中 $p\in\mathbb R$。
\SubQ{1}{$\displaystyle\sum_{n=1}^{\infty}\frac{(-1)^{n-1}}{n^{p+1/n}}$。}
\SubQ{2}{$\displaystyle\sum_{n=1}^{\infty}(-1)^n\frac{\cos2n}{n^p}$。}
\SubQ{3}{$\displaystyle\sum_{n=1}^{\infty}(-1)^{n-1}\left(e-\left(1+\frac1n\right)^n\right)$。}
\end{minipage}\par\vfill
\noindent\begin{minipage}{\textwidth}\ProblemHeading{习题}{14.107}
讨论下列级数的绝对收敛性和条件收敛性，其中 $p\in\mathbb R$。
\SubQ{1}{$\displaystyle\sum_{n=1}^{\infty}(-1)^n(n^{1/n}-1)$。}
\SubQ{2}{$\displaystyle\sum_{n=1}^{\infty}\left(a^{1/n}-\frac{b^{1/n}+c^{1/n}}2\right)$，$a,b,c>0$。}
\end{minipage}\par\vfill
\noindent\begin{minipage}{\textwidth}\ProblemHeading{习题}{14.108}
令 $a_n^+=(|a_n|+a_n)/2$，$a_n^-=(|a_n|-a_n)/2$。证明：若 $\sum a_n^+$ 与 $\sum a_n^-$ 均收敛，则 $\sum a_n$ 绝对收敛。
\end{minipage}\par\vfill
\newpage
\noindent\begin{minipage}{\textwidth}\ProblemHeading{习题}{14.109}
设 $\sum_{n=1}^{\infty}a_n$ 条件收敛，令 $a_n^+=(|a_n|+a_n)/2$，$a_n^-=(|a_n|-a_n)/2$。
\SubQ{1}{证明 $\sum a_n^+=\sum a_n^-=+\infty$，并讨论逆命题是否成立。}
\SubQ{2}{记 $S_N^+=\sum_{n=1}^Na_n^+$、$S_N^-=\sum_{n=1}^Na_n^-$。证明 $S_N^+/S_N^-\to1$。}
\end{minipage}\par\vfill
\noindent\begin{minipage}{\textwidth}\ProblemHeading{习题}{14.110}
设 $0<\alpha<1$，$p,q$ 为正整数。将 $1-1/2^\alpha+1/3^\alpha-1/4^\alpha+\cdots$ 重排为每次按原次序取 $p$ 个正项、再取 $q$ 个负项，如此反复。证明新级数收敛当且仅当 $p=q$；$p>q$ 时发散到 $+\infty$，$p<q$ 时发散到 $-\infty$。
\end{minipage}\par\vfill
\noindent\begin{minipage}{\textwidth}\ProblemHeading{习题}{14.111}
讨论下列级数的绝对收敛性和条件收敛性。
\SubQ{1}{$\displaystyle\sum_{n=2}^{\infty}\ln\left(1+\frac{(-1)^n}{n^p}\right)$，$p>0$。}
\SubQ{2}{$\displaystyle\sum_{n=1}^{\infty}(-1)^{n-1}\left(\frac{1\cdot3\cdots(2n-1)}{2\cdot4\cdots2n}\right)^p$，$p>0$。}
\SubQ{3}{$\displaystyle\frac1{1^p}-\frac1{2^q}+\frac1{3^p}-\frac1{4^q}+\frac1{5^p}-\frac1{6^q}+\cdots$，$p,q>0$。}
\end{minipage}\par\vfill
\newpage
\noindent\begin{minipage}{\textwidth}\ProblemHeading{习题}{14.112}
若对任意趋于 $0$ 的数列 $\{x_n\}$，级数 $\sum_{n=1}^{\infty}a_nx_n$ 均收敛，证明 $\sum_{n=1}^{\infty}a_n$ 绝对收敛。若只要求上述性质对每个递减趋于 $0$ 的数列成立，结论是否仍成立？
\end{minipage}\par\vfill
\noindent\begin{minipage}{\textwidth}\ProblemHeading{习题}{14.113}
只改变调和级数 $1+\frac12+\frac13+\cdots$ 各项的符号，使其依次出现 $p$ 个正项、$q$ 个负项，如此反复，而不改变分母的次序。证明新级数收敛当且仅当正整数 $p=q$。
\end{minipage}\par\vfill
\noindent\begin{minipage}{\textwidth}\ProblemHeading{习题}{14.114}
设 $\sum_{n=1}^{\infty}a_n$ 绝对收敛，$b_n\to0$。证明 $\lim_{n\to\infty}(a_1b_n+a_2b_{n-1}+\cdots+a_nb_1)=0$。
\end{minipage}\par\vfill
\newpage
\noindent\begin{minipage}{\textwidth}\ProblemHeading{例题}{14.115（改写）}
定义 $E(x)=\sum_{n=0}^{\infty}x^n/n!$。用 D’Alembert 比值判别法证明该级数对每个实数 $x$ 绝对收敛，并利用 Cauchy 乘积证明 $E(x)E(y)=E(x+y)$。
\end{minipage}\par\vfill
\noindent\begin{minipage}{\textwidth}\ProblemHeading{习题}{14.116（校订）}
设 $|q|<1$。证明\[\left(\sum_{n=0}^{\infty}q^n\right)^2=\sum_{n=0}^{\infty}(n+1)q^n.\]
\end{minipage}\par\vfill
\noindent\begin{minipage}{\textwidth}\ProblemHeading{习题}{14.117（校订）}
设幂级数 $\sum_{n=0}^{\infty}a_nx^n$、$\sum_{n=0}^{\infty}b_nx^n$ 对 $x\in(-R,R)$ 绝对收敛，$R>0$。证明\[\left(\sum_{n=0}^{\infty}a_nx^n\right)\left(\sum_{n=0}^{\infty}b_nx^n\right)=\sum_{n=0}^{\infty}c_nx^n,\qquad c_n=\sum_{k=0}^na_kb_{n-k}.\]
\end{minipage}\par\vfill
\newpage
\noindent\begin{minipage}{\textwidth}\ProblemHeading{习题}{14.118}
设\[C(x)=\sum_{n=0}^{\infty}(-1)^n\frac{x^{2n}}{(2n)!},\qquad S(x)=\sum_{n=0}^{\infty}(-1)^n\frac{x^{2n+1}}{(2n+1)!}.\]证明两级数对所有实数 $x$ 绝对收敛，且 $S(2x)=2S(x)C(x)$。
\end{minipage}\par\vfill
\noindent\begin{minipage}{\textwidth}\ProblemHeading{习题}{14.119}
设 $\alpha,\beta>0$。证明级数 $\sum_{n=1}^{\infty}(-1)^{n-1}/n^\alpha$ 与 $\sum_{n=1}^{\infty}(-1)^{n-1}/n^\beta$ 的 Cauchy 乘积在 $\alpha+\beta>1$ 时收敛，在 $\alpha+\beta\le1$ 时发散。
\end{minipage}\par\vfill
\noindent\begin{minipage}{\textwidth}\ProblemHeading{习题}{14.120}
证明以下两个发散级数的 Cauchy 乘积绝对收敛：\[1-\sum_{n=1}^{\infty}\left(\frac32\right)^n,\qquad1+\sum_{n=1}^{\infty}\left(\frac32\right)^{n-1}\left(2^n+\frac1{2^{n+1}}\right).\]
\end{minipage}\par\vfill
\newpage
\noindent\begin{minipage}{\textwidth}\ProblemHeading{例题}{14.121（改写）}
证明无穷乘积 $\prod_{n=2}^{\infty}(1-1/n^2)$ 收敛，并求其值（无穷乘积收敛指部分乘积有非零有限极限）。
\end{minipage}\par\vfill
\noindent\begin{minipage}{\textwidth}\ProblemHeading{例题}{14.122}
对实数 $x\ne0$，证明\[\lim_{N\to\infty}\prod_{n=1}^N\cos\frac{x}{2^n}=\frac{\sin x}{x}.\]当右侧为零时，此式表示部分乘积极限为零。
\end{minipage}\par\vfill
\noindent\begin{minipage}{\textwidth}\ProblemHeading{例题}{14.123}
证明 Wallis 公式 $\displaystyle\prod_{k=1}^{\infty}\frac{(2k)^2}{(2k-1)(2k+1)}=\frac\pi2$。
\end{minipage}\par\vfill
\newpage
\noindent\begin{minipage}{\textwidth}\ProblemHeading{例题}{14.124（改写）}
令 $p_n=1/2$（$n\ge1$）。计算 $\prod_{k=1}^np_k$ 的极限，并说明按部分乘积必须趋于非零有限值的定义，$\prod_{n=1}^{\infty}p_n$ 发散。
\end{minipage}\par\vfill
\noindent\begin{minipage}{\textwidth}\ProblemHeading{例题}{14.125（改写）}
证明以下部分乘积极限：
\SubQ{1}{$\displaystyle\lim_{N\to\infty}\prod_{n=1}^N\left(1+\frac1n\right)=+\infty$。}
\SubQ{2}{$\displaystyle\lim_{N\to\infty}\prod_{n=2}^N\left(1-\frac1n\right)=0$。}
\end{minipage}\par\vfill
\noindent\begin{minipage}{\textwidth}\ProblemHeading{例题}{14.126（改写）}
证明无穷乘积 $\displaystyle\prod_{n=2}^{\infty}\left(1+\frac1{n^2}\right)$ 与 $\displaystyle\prod_{n=2}^{\infty}\left(1-\frac1{n^2}\right)$ 均收敛（部分乘积极限须为非零有限值）。
\end{minipage}\par\vfill
\newpage
\noindent\begin{minipage}{\textwidth}\ProblemHeading{例题}{14.127}
设 $\alpha>-1$，$\binom\alpha0=1$，$\binom\alpha n=\alpha(\alpha-1)\cdots(\alpha-n+1)/n!$。证明 $\lim_{n\to\infty}\binom\alpha n=0$。
\end{minipage}\par\vfill
\noindent\begin{minipage}{\textwidth}\ProblemHeading{例题}{14.128（改写）}
设 $\alpha>0$。讨论无穷乘积 $\displaystyle\prod_{n=1}^{\infty}\left(1+\frac{(-1)^{n-1}}{n^\alpha}\right)$ 的敛散性。其中部分乘积有非零有限极限时称乘积收敛；若收敛且 $\displaystyle\prod_{n=1}^{\infty}(1+n^{-\alpha})$ 也收敛，称原乘积绝对收敛，否则称其条件收敛。
\SubQ{1}{确定原无穷乘积何时收敛，何时部分乘积趋于零。}
\SubQ{2}{当 $\alpha=1$ 时，判断原无穷乘积是绝对收敛还是条件收敛。}
\end{minipage}\par\vfill
\noindent\begin{minipage}{\textwidth}\ProblemHeading{习题}{14.129}
已知无穷乘积 $\prod_{n=1}^{\infty}p_n$ 与 $\prod_{n=1}^{\infty}q_n$ 的敛散性，能否据此判断下列乘积的敛散性？收敛均指部分乘积极限非零且有限。
\SubQ{1}{$\displaystyle\prod_{n=1}^{\infty}(p_n+q_n)$。}
\SubQ{2}{$\displaystyle\prod_{n=1}^{\infty}p_nq_n$。}
\SubQ{3}{$\displaystyle\prod_{n=1}^{\infty}\frac{p_n}{q_n}$，假设 $q_n\ne0$。}
\end{minipage}\par\vfill
\newpage
\noindent\begin{minipage}{\textwidth}\ProblemHeading{习题}{14.130（校订）}
证明下列等式。
\SubQ{1}{$\displaystyle\prod_{n=2}^{\infty}\frac{n^3-1}{n^3+1}=\frac23$。}
\SubQ{2}{$\displaystyle\prod_{n=2}^{\infty}\left(1-\frac2{n(n+1)}\right)=\frac13$。}
\SubQ{3}{$\displaystyle\prod_{n=0}^{\infty}\left(1+\left(\frac12\right)^{2^n}\right)=2$。}
\end{minipage}\par\vfill
\noindent\begin{minipage}{\textwidth}\ProblemHeading{习题}{14.131（校订）}
利用 Wallis 公式证明：
\SubQ{1}{$\displaystyle\prod_{n=1}^{\infty}\left(1-\frac1{4n^2}\right)=\frac2\pi$。}
\SubQ{2}{$\displaystyle\prod_{n=1}^{\infty}\left(1-\frac1{(2n+1)^2}\right)=\frac\pi4$。}
\end{minipage}\par\vfill
\noindent\begin{minipage}{\textwidth}\ProblemHeading{习题}{14.132}
讨论下列无穷乘积的敛散性（收敛须有非零有限的部分乘积极限）。
\SubQ{1}{$\displaystyle\prod_{n=1}^{\infty}\frac1n$。}
\SubQ{2}{$\displaystyle\prod_{n=1}^{\infty}\frac{(n+1)^2}{n(n+2)}$。}
\SubQ{3}{$\displaystyle\prod_{n=1}^{\infty}\sqrt[n]{1+\frac1n}$。}
\end{minipage}\par\vfill
\newpage
\noindent\begin{minipage}{\textwidth}\ProblemHeading{习题}{14.133}
讨论下列无穷乘积的敛散性（收敛须有非零有限的部分乘积极限）。
\SubQ{1}{$\displaystyle\prod_{n=1}^{\infty}\frac n{\sqrt{n^2+1}}$。}
\SubQ{2}{$\displaystyle\prod_{n=2}^{\infty}\left(\frac{n^2-1}{n^2+1}\right)^p$，$p\in\mathbb R$。}
\SubQ{3}{$\displaystyle\prod_{n=1}^{\infty}\sqrt[n]{\ln(n+x)-\ln n}$，$x>0$。}
\end{minipage}\par\vfill
\noindent\begin{minipage}{\textwidth}\ProblemHeading{习题}{14.134}
设 $\sum_{n=1}^{\infty}a_n^2<+\infty$，且每个 $\cos a_n\ne0$。证明无穷乘积 $\prod_{n=1}^{\infty}\cos a_n$ 收敛。
\end{minipage}\par\vfill
\noindent\begin{minipage}{\textwidth}\ProblemHeading{习题}{14.135}
设 $a_n>-1$，无穷乘积 $\prod(1+a_n)$ 收敛，而 $\prod(1+|a_n|)$ 发散。证明可适当重排因子使部分乘积趋于任意预先给定的正数 $s$，也可以使其趋于 $+\infty$ 或 $0$。
\end{minipage}\par\vfill
\newpage
\noindent\begin{minipage}{\textwidth}\ProblemHeading{习题}{14.136}
证明\[\frac2{\sqrt2}\cdot\frac2{\sqrt{2+\sqrt2}}\cdot\frac2{\sqrt{2+\sqrt{2+\sqrt2}}}\cdots=\frac\pi2.\]
\end{minipage}\par\vfill
\noindent\begin{minipage}{\textwidth}\ProblemHeading{习题}{14.137}
设 $a_1=1$，$a_n=n(a_{n-1}+1)$（$n\ge2$）。证明 $\prod_{n=1}^{\infty}(1+1/a_n)=e$。
\end{minipage}\par\vfill
\noindent\begin{minipage}{\textwidth}\ProblemHeading{习题}{14.138（校订）}
对 $k\ge2$ 定义 $a_{2k-1}=-1/\sqrt k$，$a_{2k}=1/\sqrt k+1/k+1/(k\sqrt k)$。证明 $\sum_{n=3}^{\infty}a_n$ 与 $\sum_{n=3}^{\infty}a_n^2$ 均发散，而 $\prod_{n=3}^{\infty}(1+a_n)$ 收敛。
\end{minipage}\par\vfill
\newpage
\noindent\begin{minipage}{\textwidth}\ProblemHeading{习题}{14.139}
证明\[\lim_{N\to\infty}\prod_{n=1}^N\frac1e\left(1+\frac1n\right)^n=0.\]由此证明 $\lim_{n\to\infty}(n^n/n!)e^{-n}=0$。
\end{minipage}\par\vfill
\noindent\begin{minipage}{\textwidth}\ProblemHeading{习题}{14.140}
若级数 $\sum_{n=1}^{\infty}a_n(x^2-1)(x^2-2^2)\cdots(x^2-n^2)$ 在某个非整数 $x=x_0$ 处收敛，证明它对所有实数 $x$ 都收敛。
\end{minipage}\par\vfill
\noindent\begin{minipage}{\textwidth}\ProblemHeading{习题}{14.141}
设 $a_n>0$，且\[\frac{a_n}{a_{n+1}}=1+\frac1n+O(b_n)\quad(n\to\infty),\]其中 $\sum_{n=1}^{\infty}b_n$ 绝对收敛。证明 $\sum_{n=1}^{\infty}a_n$ 发散。
\end{minipage}\par\vfill
\newpage
\noindent\begin{minipage}{\textwidth}\ProblemHeading{习题}{14.142}
若正项级数 $\sum_{n=1}^{\infty}a_n$ 发散，证明\[\sum_{n=1}^{\infty}\frac{a_n}{(1+a_1)(1+a_2)\cdots(1+a_n)}=1.\]
\end{minipage}\par\vfill
